\documentclass[10pt,a4paper]{amsart}
\usepackage[margin=19mm]{geometry}
\usepackage{amsmath,amssymb,mathtools}
\usepackage[scr=rsfs,cal=euler]{mathalfa}
\usepackage{xfrac}
\usepackage[unicode,hidelinks,bookmarks=false]{hyperref}
\newcommand{\ie}{i.e.\ }
\newcommand{\cf}{cf.\ }
\newcommand{\resp}{respectively\ }
\newcommand{\Subsection}{\subsection}
\providecommand{\nobreakdash}{\nobreak\hbox{-}}
\setlength{\emergencystretch}{2em}
\multlinegap0pt
\usepackage{amssymb}
\usepackage{mathtools}
\usepackage{dsfont}
\usepackage{bm}
\usepackage{braket}
\usepackage{tikz}
\usepackage{xcolor}
\newtheorem{lemma}{Lemma}
\title{Quantum Gibbs Sampling in infinite dimensions:\\ \tiny{Generation, mixing times and circuit implementation}}
\author{Simon Becker, Cambyse Rouz\'e, Robert Salzmann}
\date{Source: arXiv:2604.01192v1 (CC BY 4.0)}
\begin{document}
\maketitle

\noindent\emph{Source and licence.} Quantum Gibbs Sampling in infinite dimensions:\\ \tiny{Generation, mixing times and circuit implementation}. Version arXiv:2604.01192v1 (CC BY 4.0), distributed by arXiv under the Creative Commons Attribution 4.0 International licence, http://creativecommons.org/licenses/by/4.0/, as recorded in the arXiv licence metadata for this exact version and shown on the abstract page https://arxiv.org/abs/2604.01192v1. Full licence notice: \texttt{licenses/arxiv-2604.01192-CC-BY-4.0.txt}.\par\smallskip

\noindent\textit{Editorial scope.} Counted unit: the article's lemma lem:Riesz\_finite\_lattices (source0.tex lines 1115--1136). The article's complete original proof of it is delivered with it (source0.tex lines 1138--1319, one \texttt{proof} environment of 182 lines). No mathematics is written, repaired or re-derived: the statement and the proof are the article's own lines, in the article's own notation, and no retained line is substituted or re-flowed.  No further source-local context is needed: every cross-reference of every retained line resolves inside the two delivered spans.  Nothing was omitted from any retained span, so the package carries no omission record.  No retained line contains a citation, so the wrapper carries no bibliography and no bibliography entry.  No retained line contains a figure environment, an image-inclusion command or drawing code, so no image is delivered and the package declares no asset dependency.  Editorial formatting only, in addition: an AMS \texttt{amsart} preamble that restates the article's own declarations and macros used by the retained lines, namely the environments \texttt{lemma} on separate counters, together with the standard packages those lines need.  The retained environments print in this document as Lemma 1 in order of appearance, whereas the article prints the same items with the numbering of its own full document.  The complete native source of the article as distributed in the arXiv e-print for this exact version is bundled unchanged as \texttt{sources/197642/source0.tex}.
\begin{lemma}[Stable phase retrieval]\label{lem:Riesz_finite_lattices}
Let
\[
w(t):=\frac{1}{\beta\cosh(2\pi t/\beta)}, 
\qquad
k(s):=\int_{\mathbb R} w(t)e^{its}\,dt
= \frac{1}{2\cosh(\beta s/4)}.
\]
Let $\omega>0$ and assume that the frequency set $\Lambda\subset\mathbb R$ satisfies
\[
\Lambda \subset \bigcup_{r=1}^M (a_r+\omega\mathbb Z),
\]
for some $M\in\mathbb N$ and pairwise distinct $a_1,\dots,a_M\in[0,\omega)$.

Then there exists a constant $A>0$ such that for every finitely supported family $(c_\nu)_{\nu\in\Lambda}$ in a Hilbert space $\mathcal H$,
\begin{equation}\label{eq:Riesz_lower}
A \sum_{\nu\in\Lambda} \|c_\nu\|_{\mathcal H}^2
\;\le\;
\int_{\mathbb R} w(t)\Big\|\sum_{\nu\in\Lambda} c_\nu e^{it\nu}\Big\|_{\mathcal H}^2\,dt.
\end{equation}
In addition, if the non-zero Bohr frequencies are at least $\delta>0$, then for $\beta>0$ large enough, then $\mathcal E(x)=0$ implies that $x_{\nu}=0.$
\end{lemma}

\begin{proof}
For each $r=1,\dots,M$, set
\[
\Lambda_r := \Lambda \cap (a_r+\omega\mathbb Z),
\qquad
c_{r,n} := c_{a_r+n\omega},
\]
with the convention that $c_{r,n}=0$ if $a_r+n\omega\notin\Lambda$.
Then
\[
\sum_{\nu\in\Lambda} c_\nu e^{it\nu}
=
\sum_{r=1}^M e^{ita_r}\sum_{n\in\mathbb Z} c_{r,n} e^{in\omega t}.
\]
Define
\[
F_r(t):=\sum_{n\in\mathbb Z} c_{r,n} e^{in\omega t},
\qquad
F(t):=\sum_{r=1}^M e^{ita_r}F_r(t).
\]

Expanding the square and using the definition of $k$, we obtain
\begin{align*}
\int_{\mathbb R} w(t)\|F(t)\|_{\mathcal H}^2\,dt
&=
\sum_{r,s=1}^M \sum_{n,m\in\mathbb Z}
\langle c_{r,n},c_{s,m}\rangle_{\mathcal H}
\int_{\mathbb R} w(t)e^{it[(a_r-a_s)+(n-m)\omega]}\,dt \\
&=
\sum_{r,s=1}^M \sum_{n,m\in\mathbb Z}
k\big((a_r-a_s)+(n-m)\omega\big)\,
\langle c_{r,n},c_{s,m}\rangle_{\mathcal H}.
\end{align*}

For $\theta\in[-\pi,\pi]$, define
\[
\widehat c_r(\theta):=\sum_{n\in\mathbb Z} c_{r,n}e^{in\theta},
\]
and define the scalar symbols
\[
m_{rs}(\theta)
:=
\sum_{\ell\in\mathbb Z}
k\big((a_r-a_s)+\ell\omega\big)e^{-i\ell\theta}.
\]
Since $k$ decays exponentially, the series defining $m_{rs}$ converges absolutely and uniformly, hence $m_{rs}$ is continuous. Moreover,
\[
m_{sr}(\theta)=\overline{m_{rs}(\theta)},
\]
so $M(\theta):=(m_{rs}(\theta))_{r,s=1}^M$ is self-adjoint.

By setting $\ell=n-m$, we obtain
\[
\sum_{n,m\in\mathbb Z}
k\bigl((a_r-a_s)+(n-m)\omega\bigr)\,
\langle c_{r,n},c_{s,m}\rangle_{\mathcal H}
= \sum_{\ell\in\mathbb Z}
k\bigl((a_r-a_s)+\ell\omega\bigr)
\sum_{m\in\mathbb Z}
\langle c_{r,m+\ell},c_{s,m}\rangle_{\mathcal H}.
\]
By Plancherel on $\ell^2(\mathbb Z)$,
\[
\sum_{m\in\mathbb Z}
\langle c_{r,m+\ell},c_{s,m}\rangle_{\mathcal H}
=
\int_{-\pi}^{\pi}
e^{-i\ell\theta}\,
\langle \widehat c_r(\theta),\widehat c_s(\theta)\rangle_{\mathcal H}
\,\frac{d\theta}{2\pi}.
\]
Hence
\[
\sum_{\ell\in\mathbb Z}
k\bigl((a_r-a_s)+\ell\omega\bigr)
\sum_{m\in\mathbb Z}
\langle c_{r,m+\ell},c_{s,m}\rangle_{\mathcal H}= \int_{-\pi}^{\pi}
\Bigl(
\sum_{\ell\in\mathbb Z}
k\bigl((a_r-a_s)+\ell\omega\bigr)e^{-i\ell\theta}
\Bigr)
\langle \widehat c_r(\theta),\widehat c_s(\theta)\rangle_{\mathcal H}
\frac{d\theta}{2\pi}.
\]
This means that
\[
\int_{\mathbb R} w(t)\|F(t)\|_{\mathcal H}^2\,dt
=
\int_{-\pi}^{\pi}
\sum_{r,s=1}^M
m_{rs}(\theta)\,
\langle \widehat c_r(\theta),\widehat c_s(\theta)\rangle_{\mathcal H}
\,\frac{d\theta}{2\pi}.
\]

We now show that $M(\theta)$ is strictly positive definite for each $\theta$.
By Poisson summation,
\[
m_{rs}(\theta)
=
\frac{2\pi}{\omega}\sum_{j\in\mathbb Z}
w\!\left(\frac{\theta+2\pi j}{\omega}\right)
e^{\,i(a_r-a_s)(\theta+2\pi j)/\omega}.
\]
Hence, for any $z=(z_1,\dots,z_M)\in\mathbb C^M$,
\begin{align*}
\sum_{r,s=1}^M m_{rs}(\theta) z_r\overline{z_s}
&=
\frac{2\pi}{\omega}\sum_{j\in\mathbb Z}
w\!\left(\frac{\theta+2\pi j}{\omega}\right)
\sum_{r,s=1}^M
z_r\overline{z_s}
e^{\,i(a_r-a_s)(\theta+2\pi j)/\omega} \\
&=
\frac{2\pi}{\omega}\sum_{j\in\mathbb Z}
w\!\left(\frac{\theta+2\pi j}{\omega}\right)
\Big|\sum_{r=1}^M z_r e^{ia_r(\theta+2\pi j)/\omega}\Big|^2.
\end{align*}
This shows that $M(\theta)$ is positive semidefinite.

Assume now that
\[
\sum_{r,s=1}^M m_{rs}(\theta) z_r\overline{z_s}=0.
\]
Since $w>0$, every term in the above sum must vanish, hence
\[
\sum_{r=1}^M z_r e^{ia_r(\theta+2\pi j)/\omega}=0
\qquad\text{for every }j\in\mathbb Z.
\]
Set
\[
\lambda_r:=e^{2\pi i a_r/\omega},
\qquad
c_r:=z_r e^{ia_r\theta/\omega}.
\]
Then the preceding identities become
\[
\sum_{r=1}^M c_r \lambda_r^j = 0
\qquad\text{for every }j\in\mathbb Z.
\]
Since the $a_r$ are distinct modulo $\omega$, the numbers $\lambda_r$ are pairwise distinct. Taking $j=0,\dots,M-1$, we obtain a Vandermonde system, hence
\[
c_r=0 \quad\text{for all }r.
\]
Therefore $z_r=0$ for all $r$, and so $M(\theta)$ is strictly positive definite.

Let $\lambda_{\min}(\theta)$ denote the smallest eigenvalue of $M(\theta)$. Since $M(\theta)$ depends continuously on $\theta$ and is strictly positive definite for every $\theta$, compactness of $[-\pi,\pi]$ yields
\[
A:=\min_{\theta\in[-\pi,\pi]} \lambda_{\min}(\theta)>0.
\]
Therefore
\[
\sum_{r,s=1}^M
m_{rs}(\theta)\,
\langle u_r,u_s\rangle_{\mathcal H}
\ge
A\sum_{r=1}^M \|u_r\|_{\mathcal H}^2
\qquad
\text{for all }u_1,\dots,u_M\in\mathcal H.
\]
Applying this with $u_r=\widehat c_r(\theta)$ and integrating, we get
\[
\int_{\mathbb R} w(t)\|F(t)\|_{\mathcal H}^2\,dt
\ge
A\int_{-\pi}^{\pi}\sum_{r=1}^M \|\widehat c_r(\theta)\|_{\mathcal H}^2\,\frac{d\theta}{2\pi}.
\]
Finally, Parseval gives
\[
\int_{-\pi}^{\pi}\|\widehat c_r(\theta)\|_{\mathcal H}^2\,\frac{d\theta}{2\pi}
=
\sum_{n\in\mathbb Z}\|c_{r,n}\|_{\mathcal H}^2,
\]
and so
\[
\int_{\mathbb R} w(t)\|F(t)\|_{\mathcal H}^2\,dt
\ge
A\sum_{r=1}^M\sum_{n\in\mathbb Z}\|c_{r,n}\|_{\mathcal H}^2
=
A\sum_{\nu\in\Lambda}\|c_\nu\|_{\mathcal H}^2.
\]
This proves \eqref{eq:Riesz_lower}.
\end{proof}

\end{document}
